PID controllers still regulate the large majority of loops in the chemical process industries \cite{astrom2006pid, seborg2010process}, and the temperature loop of an exothermic continuous stirred tank reactor (CSTR) is one of the most demanding among them. Arrhenius kinetics couple reactor temperature and reaction rate through an exponential positive feedback, which produces steady-state multiplicity, parametric sensitivity and open-loop unstable regimes \cite{uppal1974dynamic, aris1958analysis, bequette2003process}. The cooling valve that closes the loop is bounded in stroke and in travel speed, and when an aggressive controller drives it into saturation the loop is effectively opened and the integrator winds up \cite{hippe2006windup, tarbouriech2009antiwindup}.

The tuning rules most often used as references for this loop are the reaction-curve rules of Ziegler and Nichols (1942) \cite{ziegler1942optimum} and Cohen and Coon (1953) \cite{cohen1953theoretical}, derived for linear, unconstrained processes \cite{odwyer2009handbook, shinskey1996process}. Their failure on exothermic reactors is a recurring finding \cite{visioli2006practical, luyben2007chemical}, and it is frequently used to motivate model-based, nonlinear or learning-based alternatives. Two aspects of such comparisons are rarely examined. First, the empirical rules are a strawman: modern robust rules such as Tyreus--Luyben \cite{tyreus1992tuning}, AMIGO \cite{astrom2004revisiting}, Skogestad's SIMC \cite{skogestad2003simple, grimholt2018optimal} and $M_s$-constrained optimization \cite{garpinger2014performance, padula2012optimal} are the relevant linear baselines. Second, the verdict on a tuning rule depends on how the actuator and the anti-windup logic are modelled. A clamping rule that checks only the amplitude limits, and that ignores the sign of a reverse-acting loop, keeps integrating while a rate-limited valve ramps towards its target; the resulting windup is then attributed to the tuning rule rather than to the implementation.

This paper separates these effects on a three-state CSTR model with an explicit jacket balance and a cooling valve limited to $[0, 300]$~L/min and $80$~L/min$^2$. Eight PID tunings are compared under one protocol, with error metrics evaluated on the true reactor temperature. The contributions are:
\begin{enumerate}
    \item A directional clamping anti-windup rule that freezes the integrator whenever the previous command was cut by the amplitude \emph{or} the slew-rate limit and the integral increment pushes further into saturation, with the sign of the reverse-acting loop built in. An ablation against the textbook amplitude-only rule shows that the implementation, not the tuning rule, decides whether the empirical rules lose the reactor.
    \item A benchmark of eight tunings---Ziegler--Nichols, Cohen--Coon, Tyreus--Luyben, AMIGO and its two-degree-of-freedom (2-DOF) variant, SIMC, and multi-start ITAE optimization with $M_s \le 1.6$ evaluated on the identified dead-time loop, with its 2-DOF variant---covering servo tracking, unmeasured load disturbances and the robustness margins of the design model and of the exact linearized plant.
    \item A continuous map of the trade-off between tracking error and valve travel obtained by sweeping the SIMC closed-loop time constant and by scalarized optimization \cite{alcantara2013simple, reynoso2021multiobjective}.
    \item A stochastic and global sensitivity analysis: $1200$ Monte Carlo runs under thermocouple noise and Sobol' indices \cite{sobol2001global, saltelli2010variance} of the tracking error and valve travel under simultaneous uncertainty in kinetics, heat transfer, feed concentration and sensor dead time.
\end{enumerate}
Section~\ref{sec:model} presents the plant and the actuator limits, Section~\ref{sec:pid} the controller structure, the anti-windup rule and the tuning rules, Section~\ref{sec:results} the results and Section~\ref{sec:conclusion} the conclusions.
